\section{Thesis claim and research questions}\label{sec:introduction:research_questions}
We do not need to understand how electromagnetism propagates through a transistor to program a classical computer.
As Núñez-Corrales et al.~observe of mature classical programming systems:

\begin{quote}
    \textit{When reasoning about the effects of a program, a programmer works at the level of formal statements (mathematics), translates statements into programs (instructions) and ideally uses knowledge of the limitations of finite representations to validate program correctness (proofs).
    Programmers no longer need to worry about the fate of individual bits or the sequence of gates the machine executes~\cite{nunez_corrales2025betterabstractmachines}.}
\end{quote}

There is no corresponding principle in quantum programming.
Programmers must currently understand how quantum mechanics operates at the gate and circuit level to write correct programs.
This thesis introduces and investigates a class of language that extends the classical principle to the quantum domain.
A \textit{quantum-implicit} programming language is one whose primary surface model is expressed in ordinary mathematical and computational abstractions.
It allows programmers to describe and implement quantum algorithms without explicitly annotating or reasoning about qubits, gates, circuits or other quantum-mechanical constructs.
The compiler is responsible for inferring and synthesising the necessary quantum realisation.

\subsection{Claim}\label{subsec:introduction:research_questions:claim}
In a quantum-implicit language, programming looks substantively
different~\cite{nunez_corrales2025betterabstractmachines}:
\begin{itemize}
    \item Programs are expressed at a level of algorithmic intent rather than circuit structure;
    \item Classical and quantum data types interoperate without the programmer managing qubit modality or gate sets;
    \item The number of execution shots required to produce a statistically reliable distribution is determined by the compiler or runtime, not the programmer; and
    \item Complexity advantages over classical computation are assessable systematically,without hardware-specific reasoning.
\end{itemize}

Table~\ref{tab:introduction:research_questions:criteria_violations} illustrates why this tier is currently unoccupied.
Three of the most advanced existing high-level languages are assessed against four criteria governing the programming layer~\cite{nunez_corrales2025betterabstractmachines}; each records at least one violation.

\begin{table}[ht]
    \centering
    \caption{Representative criteria violations across three advanced high-level quantum programming languages. Each language records at least one violation of the four criteria most directly governing the programming layer~\cite{nunez_corrales2025betterabstractmachines}.
    The full criteria set is presented in Chapter~\ref{ch:literature_review}.}
    \label{tab:introduction:research_questions:criteria_violations}
    \small
    \begin{tabular}{p{3.2cm} p{3.5cm} p{3.5cm} p{3.5cm}}
        \toprule
        \textbf{Criterion} & \textbf{Silq} & \textbf{Quff} & \textbf{Qutes} \\\midrule
        \textbf{Symbolic denotational syntax.}
        No circuit elements required to reason about program effects on state.
        & \texttt{qfree}, \texttt{const}, \texttt{lifted} and \texttt{mfree} annotations encode circuit-execution semantics as mandatory surface syntax. Both require circuit-model knowledge to apply correctly.
        & Runtime function inversion (\texttt{invertF}) is a circuit-level surface concept. Phase notation and explicit qubit operations remain at the programming surface.
        & \texttt{ref} advances toward quantum-agnostic pass-by-reference. Explicit gate instructions (\texttt{hadamard}, \texttt{mcp}, \texttt{swap}) are still required for algorithms outside the built-in set. \\ \addlinespace
        \textbf{Representation-independent data types.}
        No reference to circuit elements in data types.
        & \texttt{uint[n]} hard-codes qubit count \texttt{n} into every numeric type. Every declaration couples data representation to a physical resource quantity.
        & Explicit bit-width specification is required for all quantum types. The programmer must declare and manage qubit count for each variable.
        & Quantum register types carry explicit size or qubit-level structure for operations outside the built-in set. \\ \addlinespace
        \textbf{Classical-quantum regularity.}
        Classical and quantum instructions coexist at the same level.
        & The \texttt{!} annotation mandates a programmer-visible stratification between classical and quantum values, with different rules for control flow and function application.
        & Distinct quantum and classical type families are required. Programmers must reason about which operations are valid in each context.
        & Classical and quantum contexts remain syntactically distinguished. The classical/quantum boundary is programmer-managed rather than compiler-inferred. \\ \addlinespace
        \textbf{Intrinsic ensemble semantics.}
        Shot count elided; distribution types well-defined at the language level.
        & Explicit \texttt{measure} is required to obtain classical output. Shot management is external to the language.
        & Measurement must be explicitly inserted by the programmer. Shot count for statistical reliability is managed outside the language.
        & Scope-bounded uncomputation handles one class of quantum-to-classical transition. General measurement and shot management remain explicit. \\
        \bottomrule
    \end{tabular}
\end{table}

Violations are consistent across the three most advanced high-level languages in the reviewed literature, confirming that the gap is a property of the existing design space rather than a shortcoming of any single system.
The full criteria set and a systematic assessment across all reviewed languages is presented in Chapter~\ref{ch:literature_review}.

\subsection{Research questions}\label{subsec:introduction:research_questions:questions}
The central research question (CRQ) this thesis addresses is:

\begin{quote}
    \textit{\textbf{CRQ: }Can a quantum-implicit programming language, whose surface model is agnostic of the underlying physical quantum substrate, be formally designed, proven correct and compiled to execute on gate-based quantum hardware?}
\end{quote}

This thesis answers that question affirmatively across four incorporated publications.

A quantum-implicit programming language can be formally specified, proven type-safe and compiled to execute on gate-based quantum hardware with proven semantic preservation.
The quantum-implicit tier is not merely aspirational; it is formally achievable.

The CRQ is decomposed into three subsidiary research questions (SRQs), each addressed by one or more thesis chapters:

\begin{quote}
    \textit{\textbf{SRQ1: }Can a quantum-implicit language be formally specified through computational abstractions requiring no knowledge of quantum mechanical principles or the circuit execution model, while remaining expressive enough to encode quantum algorithms?}
\end{quote}

SRQ1 is addressed in Chapter~\ref{ch:pub_language_design}, which presents the syntax,operational semantics and type system of the quantum-implicit language.

\begin{quote}
    \textit{\textbf{SRQ2: }Can the core source language be proven type-safe: specifically, can \emph{progress} (well-typed programs never reach a stuck state) and \emph{preservation} (evaluation preserves the type of well-typed expressions) be formally established for the language's type system?}
\end{quote}

SRQ2 is addressed in Chapter~\ref{ch:pub_formal_theory}, which presents the progress and preservation proofs for the core language.

\begin{quote}
    \textit{\textbf{SRQ3: }Can semantic preservation be formally established for the compilation pipeline: specifically, does the compiler's translation from source language to quantum IR faithfully preserve the meaning of source programs?}
\end{quote}

SRQ3 is addressed in Chapters~\ref{ch:pub_formal_theory} and~\ref{ch:pub_compilation}: the former establishes the formal semantic preservation proof; the latter demonstrates it end-to-end for concrete programs compiled to and executed on NISQ hardware.

The CRQ is answered collectively in Chapters~\ref{ch:results_evaluation} and~\ref{ch:study_implications}, which synthesise the contributions of all four publications.